% Generated by tex/build_swebok_structure.py from tex/chapters/ch01/04_要求獲得.tex
\subsubsection{要求の源泉}

要求は多くの源泉から来る。人に由来する源泉と、人以外に由来する源泉の両方を考える必要がある。

\par\smallskip\noindent\refstepcounter{table}\label{tab:ch01-06}\textbf{表 \thetable: 第01章 ソフトウェア要求 表06}\par\smallskip
\begin{tabularx}{0.97\linewidth}{@{}p{.27\linewidth}X@{}}
\toprule
源泉 & 例 \\
\midrule
\term{顧客} & ソフトウェア開発に費用を出す人・組織。 \\
\term{利用者} & ソフトウェアを直接または間接に使う人。利用頻度、権限、熟練度で分類できる。 \\
\term{業務専門家} & 対象業務、規制、現場作業に詳しい人。 \\
\term{運用担当} & 本番環境、監視、障害対応、バックアップなどに責任を持つ人。 \\
\term{サポート担当} & 利用者からの問い合わせや障害報告を受ける人。 \\
\term{規制当局・専門団体} & 法令、標準、ガイドラインを通じて要求を課す組織。 \\
\term{開発者} & 技術的制約、保守性、開発効率に関する要求を示すことがある。 \\
\bottomrule
\end{tabularx}

また、要求は人からだけではなく、文書、既存システム、業務方針、計算環境、競合製品、障害履歴、標準、法規制からも得られる。

\MiniBox{利害関係者分析}{利害関係者分析とは、関係する人や組織をできるだけ漏れなく見つけ、代表性の偏りを減らす活動である。声の大きい利用者だけから要求を集めると、少数派や運用担当の重要な要求を見落としやすい。}
